%------------------------------------------------------------------------------------
%  SECTION 4.2 : TABLEAUX AVANCÉS EN SITUATION RÉELLE
%------------------------------------------------------------------------------------
\section{Tableaux Complexes, Multipage et Orientation Paysage}
\label{sec:tableaux_avances}

La mise en forme de données expérimentales denses exige des techniques avancées pour s'adapter à la largeur de page sans tronquer le texte et sans générer d'erreurs de dépassement (\texttt{overfull \textbackslash hbox}).

\subsection{Dimensionnement automatique avec \textsf{tabularx}}
\label{subsec:tabularx_auto}

En \LaTeX{} standard, l'environnement \texttt{tabular} requiert de fixer manuellement la largeur de chaque colonne (ex. \texttt{p\{4cm\}}).
Si le texte est trop long, il déborde de la marge de la page.
À l'inverse, l'environnement \texttt{tabularx} introduit le spécificateur de colonne majuscule \textbf{\texttt{X}} :
\begin{itemize}
  \item Il calcule automatiquement l'espace résiduel sur la largeur totale prescrite (\texttt{\textbackslash textwidth}).
  \item Il répartit proportionnellement cet espace entre toutes les colonnes \texttt{X}.
  \item Il justifie automatiquement le texte à l'intérieur de la cellule avec césure propre.
\end{itemize}

Le~\cref{tab:comparatif_classe} met en pratique cette flexibilité pour confronter les atouts de \texttt{ClasseGuillaume} face à une installation standard :

%------------------------------------------------------------------------------------
%  TABLEAU COMPARATIF : LATEX STANDARD VS CLASSEGUILLAUME
%------------------------------------------------------------------------------------
\begin{table}[htbp]
  \centering
  \small
  \caption{Comparatif synthétique des fonctionnalités : \LaTeX{} Standard vs \texttt{ClasseGuillaume}}
  \label{tab:comparatif_classe}
  \begin{tabularx}{\textwidth}{p{3.2cm} p{5.5cm} X}
    \toprule
    \textbf{Fonctionnalité} & \textbf{En \LaTeX{} Standard (\emph{Vanilla})} & \textbf{Avec \texttt{ClasseGuillaume.cls}} \\
    \midrule
    \textbf{Gestion des polices} & Déclarations manuelles \textsf{fontspec} avec fichiers physiques \texttt{.ttf} & Détection intelligente triple : système direct, \TeX{} Live natif ou repli Overleaf. \\
    \addlinespace
    \textbf{Interligne des titres} & \texttt{\textbackslash onehalfspacing} étire affreusement les titres multi-lignes & Protection automatique : interligne 1,5 dans le corps, interligne simple strict dans les titres. \\
    \addlinespace
    \textbf{Puces de paragraphes} & \texttt{\textbackslash paragraph} renvoie à la ligne ou nécessite un réglage KOMA & Puce colorée en ligne prête à l'emploi (\texttt{\textbackslash textbullet~}) avec poursuite immédiate du texte. \\
    \addlinespace
    \textbf{Changement de marges} & Écriture manuelle lourde de \texttt{\textbackslash newgeometry} à chaque bascule & Trois commandes immédiates : \texttt{\textbackslash geometrieWord}, \texttt{\textbackslash geometrieLarge}, \texttt{\textbackslash geometrieStandard}. \\
    \addlinespace
    \textbf{Notes de révision} & \textsf{todonotes} déborde des marges, casse en UTF-8 et sans infobulle & \texttt{\textbackslash todoH} : surlignage \textsf{lua-ul}, troncature automatique > 30 car., infobulle PDF et lien retour. \\
    \addlinespace
    \textbf{Illustrations dédiées} & Schémas et photos se mélangent aux figures dans la table des figures & Flottant \texttt{illustration} avec sa propre liste indépendante (\texttt{\textbackslash listofillustrations}). \\
    \addlinespace
    \textbf{Mise en valeur visuelle} & Environnements bruts ou code \textsf{tcolorbox} complexe à réécrire & Boîtes \texttt{cadreFigure} (coins arrondis) et \texttt{boiteSynthese} (fond teinté) prêtes à l'emploi. \\
    \addlinespace
    \textbf{Logo institutionnel} & Montage artisanal de paquets d'en-tête sujet aux décalages & Couche dynamique \textsf{scrlayer} avec bascules \texttt{\textbackslash masquerLogo} et \texttt{\textbackslash afficherLogo}. \\
    \addlinespace
    \textbf{Bibliographie FR} & Chaînes en anglais (\emph{visited on}, \emph{available from}) & Localisation soignée : « consulté le », « disponible sur », liens de retour \texttt{backref}. \\
    \bottomrule
  \end{tabularx}
\end{table}



\subsection{Gestion des grands tableaux multipages avec \textsf{longtable}}
\label{subsec:longtable_multipage}

Lorsqu'un inventaire de composants ou une série temporelle de mesures s'étend sur plusieurs pages consécutives, l'environnement flottant \texttt{table} ne peut pas se scinder.
Il faut faire appel au paquet \textsf{longtable} :

\begin{boiteSynthese}{Les quatre zones clés d'un tableau \textsf{longtable}}
  Un tableau multipage professionnel doit répéter ses en-têtes pour que le lecteur ne perde jamais le fil :
  \begin{itemize}
    \item \texttt{\textbackslash endfirsthead} : En-tête principal affiché sur la toute première page.
    \item \texttt{\textbackslash endhead} : En-tête secondaire répété en haut de chaque page suivante avec la mention explicite « \textit{Suite de la page précédente} ».
    \item \texttt{\textbackslash endfoot} : Pied de tableau affiché au bas des pages intermédiaires (« \textit{Suite sur la page suivante...} »).
    \item \texttt{\textbackslash endlastfoot} : Pied de clôture définitif arborant la règle \texttt{\textbackslash bottomrule}.
  \end{itemize}
\end{boiteSynthese}

Le~\cref{tab:composants_multipage} démontre l'intégration complète de cette structure :

%------------------------------------------------------------------------------------
%  TABLEAU MULTIPAGE AVEC LONGTABLE ET BOOKTABS
%------------------------------------------------------------------------------------
\begin{small}
\begin{longtable}{p{2.5cm} p{3.2cm} p{4.5cm} p{3.5cm}}
  \caption{Inventaire exhaustif des composants expérimentaux (Tableau multipage)}
  \label{tab:composants_multipage} \\
  \toprule
  \textbf{Référence} & \textbf{Désignation} & \textbf{Spécifications Techniques} & \textbf{Statut Métrologique} \\
  \midrule
  \endfirsthead

  \multicolumn{4}{c}{\tablename\ \thetable\ -- \textit{Suite de la page précédente}} \\
  \toprule
  \textbf{Référence} & \textbf{Désignation} & \textbf{Spécifications Techniques} & \textbf{Statut Métrologique} \\
  \midrule
  \endhead

  \midrule
  \multicolumn{4}{r}{\textit{Suite sur la page suivante \dots}} \\
  \bottomrule
  \endfoot

  \bottomrule
  \endlastfoot

  CMP-101 & Capteur piézoélectrique & Plage \qtyrange{0}{50}{\kilo\pascal}, précision \qty{0.1}{\percent} & Étalonné ISO 17025 \\
  CMP-102 & Sonde thermocouple K & Plage \qtyrange{-50}{400}{\degreeCelsius}, temps de réponse \qty{20}{\milli\second} & Certifié d'usine \\
  CMP-103 & Débitmètre vortex & Diamètre intérieur \qty{25}{\milli\meter}, sortie \qtyrange{4}{20}{\milli\ampere} & Vérification annuelle \\
  CMP-104 & Manomètre différentiel & Échelle \qtyrange{0}{5}{\bar}, membrane inox 316L & En attente d'ajustage \\
  \addlinespace
  CMP-201 & Pompe volumétrique & Cylindrée \qty{12.5}{\centi\meter\cubed}, moteur brushless & Validé banc d'essai \\
  CMP-202 & Vanne proportionnelle & Régulation PID intégrée, alimentation \qty{24}{\volt} DC & Zéro dérive mesuré \\
  CMP-203 & Échangeur à plaques & Surface d'échange \qty{1.8}{\meter\squared}, flux contre-courant & Test hydrostatique OK \\
  CMP-204 & Séparateur cyclonique & Efficacité \qty{99.2}{\percent} sur particules $> \qty{5}{\micro\meter}$ & Conforme norme ATEX \\
  \addlinespace
  CMP-301 & Module d'acquisition & 16 canaux analogiques, conversion 24 bits à \qty{100}{\kilo\hertz} & Horloge synchronisée \\
  CMP-302 & Convertisseur USB-CAN & Isolation galvanique \qty{2.5}{\kilo\volt}, débit \qty{1}{\mega\bit\per\second} & Pilote validé Linux \\
  CMP-303 & Alimentation stabilisée & Tension \qtyrange{0}{30}{\volt}, courant maxi \qty{10}{\ampere} & Ondulation $< \qty{1}{\milli\volt} $ \\
  CMP-304 & Blindage électromagnétique & Atténuation \qty{65}{\decibel} entre \qty{10}{\mega\hertz} et \qty{1}{\giga\hertz} & Cage de Faraday testée \\
\end{longtable}
\end{small}



\subsection{Tableaux panoramiques en mode paysage avec \textsf{pdflscape}}
\label{subsec:mode_paysage}

Certaines synthèses transversales comportent un nombre trop élevé de paramètres pour être contenues dans la largeur portrait d'une feuille A4 (\qty{21}{\centi\meter}).
L'utilisation de l'environnement \texttt{landscape} issu du paquet \textsf{pdflscape} résout cette contrainte avec un confort de lecture total :
contrairement au paquet classique \textsf{lscape} qui laisse la page en portrait tout en tournant le texte de 90 degrés, \textsf{pdflscape} transmet une directive au visualiseur PDF pour effectuer une \textbf{véritable rotation physique de la page}.
Le document reste ainsi parfaitement lisible à l'écran sans que l'utilisateur n'ait à pivoter son moniteur.

Le~\cref{tab:evaluation_paysage} (page suivante) offre un exemple probant de matrice décisionnelle en mode paysage.

%------------------------------------------------------------------------------------
%  TABLEAU PAYSAGE LARGE (pdflscape + booktabs + tabularx)
%------------------------------------------------------------------------------------
\begin{landscape}
\begin{table}[p]
  \centering
  \small
  \caption{Matrice d'évaluation comparative multicritères des architectures logicielles et typographiques}
  \label{tab:evaluation_paysage}
  \begin{tabularx}{\linewidth}{p{3.5cm} p{3.5cm} c c c c X}
    \toprule
    \textbf{Moteur / Chaîne} & \textbf{Standard de Polices} & \textbf{UTF-8 Natif} & \textbf{Microtypo} & \textbf{Vitesse} & \textbf{Complexité} & \textbf{Conclusion Opérationnelle} \\
    \midrule
    pdf\LaTeX{} classique & T1 / Type 1 (8 bits) & Non (inputenc) & Limitée & Rapide & Élevée (bidouilles) & Obsolète pour nouveaux projets académiques. \\
    \addlinespace
    Xe\LaTeX{} & OpenType / TrueType & Oui & Basique & Moyenne & Moyenne & Préféré historiquement pour langues non latines. \\
    \addlinespace
    \textbf{Lua\LaTeX{} moderne} & \textbf{OpenType / TrueType} & \textbf{Oui (100\,\%)} & \textbf{Complète} & \textbf{Excellente} & \textbf{Épurée} & \textbf{Standard contemporain recommandé par l'UCLouvain.} \\
    \midrule
    \multicolumn{7}{l}{\textbf{Bilan comparatif des fonctionnalités de publication :}} \\
    \addlinespace
    Classe KOMA \texttt{scrbook} & Système KOMA-Script & Oui & Oui & -- & Modérée & Socle de base solide mais sans encadrements prêts à l'emploi. \\
    \addlinespace
    \texttt{ClasseGuillaume} & \textbf{Architecture Modulaire} & \textbf{Oui} & \textbf{Pénalités max} & -- & \textbf{Immédiate} & \textbf{Intègre \texttt{\textbackslash todoH}, 3 géométries, logo scrlayer et booktabs.} \\
    \bottomrule
  \end{tabularx}
  \vspace{4mm}
  \begin{minipage}{\linewidth}
    \footnotesize
    \textit{Note méthodologique :} Ce tableau panoramique démontre l'utilisation conjointe de l'environnement \texttt{landscape} fourni par le paquet \textsf{pdflscape} et de la classe de document. La rotation est physique dans le fichier PDF généré, garantissant une consultation parfaite sans devoir tourner manuellement l'écran.
  \end{minipage}
\end{table}
\end{landscape}



